% tkz-obj-eu-polygons.tex
% Copyright 2011-2026  Alain Matthes
% SPDX-License-Identifier: LPPL-1.3c
% Maintainer: Alain Matthes

\typeout{2026/01/25 5.13c tkz-obj-eu-polygons.tex}
\makeatletter
%<--------------------------------------------------------------------------–>
%                                 Polygon
%<--------------------------------------------------------------------------–>
%<----------------------------   Draw       --------------------------------–>
%
%<--------------------------------------------------------------------------–>
\def\tkzDrawPolygon{\pgfutil@ifnextchar[{\tkz@DrawPolygon}{%
                                         \tkz@DrawPolygon[]}}
\def\tkz@DrawPolygon[#1](#2,#3){%
 \begingroup
 \draw[line style,line join=round,#1] (#2)
     \foreach \pt in {#2,#3}{--(\pt)}--cycle;%
 \endgroup
}
%<--------------------------------------------------------------------------–>
\def\tkz@multipolygons#1 #2\@nil{%
\protected@edef\tkz@temp{
\noexpand \tkzDrawPolygon[\tkz@optpolygon](#1)}\tkz@temp%
\def\tkz@nextArg{#2}%
\ifx\tkzutil@empty\tkz@nextArg
     \let\next\@gobble
\fi
\next#2\@nil
}
%<--------------------------------------------------------------------------–>
\def\tkzDrawPolygons{\pgfutil@ifnextchar[{\tkz@DrawPolygons}{\tkz@DrawPolygons[]}}
\def\tkz@DrawPolygons[#1](#2){%
\xdef\tkz@optpolygon{#1}
\begingroup
   \let\next\tkz@multipolygons
   \next#2 \@nil %
\endgroup
}%
%<--------------------------------------------------------------------------–>
%<--------------------------------------------------------------------------–>
%<----------------------------   FILL       --------------------------------–>
%
%<--------------------------------------------------------------------------–>
\def\tkzFillPolygon{\pgfutil@ifnextchar[{\tkz@FillPolygon}{%
                                         \tkz@FillPolygon[]}}
\def\tkz@FillPolygon[#1](#2,#3){%
 \begingroup
     \fill[#1] (#2)
   \foreach \pt in {#2,#3}{--(\pt)}--cycle;%
\endgroup
}
%<--------------------------------------------------------------------------–>
%<--------------   Labels for Regular Polygon      -------------------------–>
%
%<--------------------------------------------------------------------------–>
\pgfkeys{/deflabelregpoly/.cd,
      name/.store in      =   \tkz@regpolname,
      sep/.store in       =   \tkz@regpolsep,
      name        = P,
      sep         = 1.1
}

\def\tkzLabelRegPolygon{\pgfutil@ifnextchar[{\tkz@labelPolygon}{%
                                             \tkz@labelPolygon[]}}
\def\tkz@labelPolygon[#1](#2)#3{%
\begingroup
\pgfqkeys{/deflabelregpoly}{#1}
 \foreach \v [count=\vi from 1]  in {#3}  {%
  \path (#2) to [pos=1.1] node {$\v$} (\tkz@regpolname\vi) ;
  \pgfnodealias{\v}{\tkz@regpolname\vi}
  }
\endgroup
}
%<----------------------------   CLIP       --------------------------------–>
%
%<--------------------------------------------------------------------------–>

\pgfkeys{tkzclippolygon/.cd,
       out code/.is if         = tkzClipOutPoly,
       out/.code               = \tkzClipOutPolyfalse}

\def\tkzClipPolygon{\pgfutil@ifnextchar[{\tkz@ClipPolygon}{\tkz@ClipPolygon[]}}
\def\tkz@ClipPolygon[#1](#2,#3){%
\tkzClipOutPolytrue
\pgfqkeys{/tkzclippolygon}{#1}%
\iftkzClipOutPoly
\clip (#2)
   \foreach \pt in {#2,#3}{--(\pt)}--cycle;
\else
   \clip(#2)
    \foreach \pt in {#2,#3}{--(\pt)}--cycle[tkzreverseclip] ;
\fi
}%
\makeatother
\endinput